\begin{corollary}[3.1.9]
\label{IV.2.3.1.9-fr}
Sean $X$ un preesquema localmente noetheriano, $\mathcal F$ un
$\mathcal O_X$-módulo cuasicoherente y $f$ una sección de $\mathcal O_X$
sobre $X$. Para que $f$ sea $\mathcal F$-regular \textup{(\textbf{0}, 15.2.2)}, es necesario y
suficiente que $\operatorname{Ass}(\mathcal F)\subset X_f$.
\end{corollary}

En efecto, es inmediato que decir que $f$ es $\mathcal F$-regular significa
que el homomorfismo canónico
$\mathcal F\to\mathcal Hom_{\mathcal O_X}(f\mathcal O_X,\mathcal F)$ es
inyectivo, y basta aplicar (3.1.8) al ideal $\mathcal J=f\mathcal O_X$.

\oldpage[IV-2]{39}
\begin{proposition}[3.1.10]
\label{IV.2.3.1.10-fr}
Sean $X$, $Y$ preesquemas localmente noetherianos y $f:X\to Y$ un morfismo
entero. Entonces, para todo $\mathcal O_X$-módulo cuasicoherente $\mathcal F$,
se tiene
$f(\operatorname{Ass}(\mathcal F))=\operatorname{Ass}(f_*(\mathcal F))$.
\end{proposition}

Como la cuestión es local sobre $Y$ y el morfismo $f$ es afín, la cuestión se reduce
inmediatamente al caso en que $Y$ es afín, es decir, al

\begin{lemma}[3.1.10.1]
\label{IV.2.3.1.10.1-fr}
Sean $A$, $B$ dos anillos noetherianos, $\rho:A\to B$ un homomorfismo de
anillos que hace de $B$ una $A$-álgebra entera, y $M$ un $B$-módulo. Entonces
los ideales primos $\mathfrak p\in\operatorname{Ass}(M_{[\rho]})$ son las
imágenes inversas por $\rho$ de los ideales primos
$\mathfrak q\in\operatorname{Ass}(M)$.
\end{lemma}

En efecto, si $\mathfrak q\in\operatorname{Ass}(M)$, $\mathfrak q$ es el
anulador en $B$ de un elemento $x\in M$, y por tanto
$\rho^{-1}(\mathfrak q)$ es el anulador en $A$ de $x$. Recíprocamente, sea
$\mathfrak p\in\operatorname{Ass}(M_{[\rho]})$, de modo que $\mathfrak p$ es
la imagen inversa por $\rho$ del anulador $\mathfrak b$ en $B$ de un elemento
$x\in M$; del primer teorema de Cohen--Seidenberg resulta entonces que existe
un ideal primo $\mathfrak q$ de $B$ que contiene a $\mathfrak b$ y cuya imagen
inversa es $\mathfrak p$ (Bourbaki, \emph{Alg. comm.}, cap.~V, \S~2, núm.~1,
cor.~2 del teo.~1); considerando uno de los ideales primos minimales entre los
contenidos en $\mathfrak q$ y que contienen a $\mathfrak b$, se puede
evidentemente suponer que $\mathfrak q$ mismo es uno de estos ideales
minimales. Pero, como $B.x\subset M$ es isomorfo a
$B/\mathfrak b$, se sabe entonces que
$\mathfrak q\in\operatorname{Ass}(B/\mathfrak b)\subset\operatorname{Ass}(M)$
(Bourbaki, \emph{Alg. comm.}, cap.~IV, \S~1, núm.~4, teo.~2).

\begin{corollary}[3.1.11]
\label{IV.2.3.1.11-fr}
Bajo las hipótesis de \textup{(3.1.10)}, para que $\mathcal F$ carezca de ciclo primo
asociado inmerso, basta que ocurra lo mismo con $f_*(\mathcal F)$.
\end{corollary}

Supongamos, en efecto, que $f_*(\mathcal F)$ no tiene ningún ciclo primo
asociado inmerso. Observemos que, si $A$ es un álgebra entera sobre un cuerpo
$k$, todos los ideales primos de $A$ son maximales (Bourbaki,
\emph{Alg. comm.}, cap.~V, \S~2, núm.~1, prop.~1); resulta entonces de
(\textbf I, 6.2.2) que las fibras de $f$ son espacios \emph{discretos}; si $x$, $x'$
son dos puntos distintos de $\operatorname{Ass}(\mathcal F)$, ninguno de ellos
puede ser adherente al otro si $f(x)=f(x')$; y, si $f(x)\ne f(x')$, (3.1.10)
y la hipótesis implican que ninguno de los dos puntos $f(x)$, $f(x')$ puede ser
adherente al otro, y por tanto lo mismo ocurre con $x$ y $x'$.

